of @@M0@@ have degrees between @@M1@@ and @@M2@@. We then proceed by induction on @@M3@@.

When @@M4@@, i.e. @@M5@@, let @@M6@@ be the common value. The canonical homomorphism (1.2.1):
@@M7@@
is then an isomorphism. Indeed, the verification is local, and by (1.2.3), we can then suppose that @@M8@@, and the result is trivial.

In the case @@M9@@, we again consider the homomorphism (1.2.1):
@@M10@@
We verify locally that this homomorphism is injective, and that its cokernel @@M11@@ is locally free of finite type, with local generators of degree between @@M12@@ and @@M13@@: indeed, in a neighborhood of every point of @@M14@@, @@M15@@ is either zero or the nonzero homogeneous component of @@M16@@ of lowest degree; using a decomposition of @@M17@@ into direct summands of the form (1.2.2), we see in the latter case, using (1.2.3), that @@M18@@ is isomorphic to the direct summand @@M19@@ of @@M20@@, whence the asserted properties. We then apply the induction hypothesis to @@M21@@, whence the desired conclusion follows at once.

\sgaStableTarget{sga6.e06.sec.1.8}1.8. Returning to the proof of Theorem 1.1, let us show that @@M22@@ is generated as a @@M23@@-module by the powers of @@M24@@.

If @@M25@@ is a @@M26@@-Module locally free of finite type, there exists, by 1.4, a graded @@M27@@-Module of finite presentation @@M28@@, flat over @@M29@@, such that @@M30@@. By 1.6, @@M31@@ admits a graded resolution by graded @@M32@@-Modules locally free of finite type. We deduce an exact sequence:
